\textbf{B) 现在让我们回到定理1的证明。}\par

首先，假设 \(u\) 可以写成形式 \(s + n\) ，其中 \(s\) 是绝对半单的，而 \(n\) 是幂零的，并且 \(s\) 与 \(n\) 可交换。设 \(\Omega\) 为 \(K\) 的一个代数闭包；则 \(u_{(\Omega)} = s_{(\Omega)} + n_{(\Omega)}\) ，其中 \(s_{(\Omega)}\) 是可对角化的，而 \(n_{(\Omega)}\) 是幂零的，并且 \(s_{(\Omega)}\) 与 \(n_{(\Omega)}\) 可交换；由引理4可知，\(s_{(\Omega)}\) 从而 \(s\) 也一样，是唯一的，多项式 \(\chi_{u_{(\Omega)}}\) 与 \(\chi_{s_{(\Omega)}}\) 在 \(\Omega[X]\) 中一致，因此多项式 \(\chi_u\) 与 \(\chi_s\) 也一致，并且 \(s\) 可以表示为 \(u\) 的、系数属于 \(\Omega\) 的多项式。这首先表明，\(u\) 的特征值与 \(s\) 的特征值相同，从而在 \(K\) 上可分（VII，第42页，命题15）；此外，由于 \(s\) 是一个 \(\Omega\) -线性组合，由 \(u\) 的幂组成，它也是同样这些幂的 \(K\) -线性组合（II，第～311页，命题～19），故存在多项式 \(P \in K[X]\) 使得 \(s = P(u)\) ，从而 \(n = Q(u)\) ，其中 \(Q(X) = X - P(X)\) 。现在让我们证明 \(Q\) （从而 \(P\) ）可以取得没有常数项。若 \(u\) 可逆，则其特征多项式有非零常数项，且凯莱-哈密顿定理（III，第～541页，命题～20）表明 1 可以表示为 \(u\) 的一个没有常数项的多项式，故断言在此情形成立。若 \(u\) 不可逆，则其核 \(W\) 非零且在 \(n\) 下封闭（VII，第41页，引理3）；由于 \(n\) 在 \(W\) 上的限制是幂零的，故存在向量 \(x \neq 0\) 属于 \(W\) ，使得 \(u(x) = n(x) = 0\) ，这表明 \(Q\) 不能有常数项。

反之，假设 \(u\) 的特征值在 \(K\) 上可分，并设 \(L\) 是包含这些特征值的、\(K\) 的有限伽罗瓦扩张。根据引理4，我们可以写出 \(u_{(L)} = v + w\) ，其中 \(v\) 是可对角化的，而 \(w\) 是幂零的，并且 \(vw = wv\) 。设 \(B\) 为 \(E\) 的一组基，设 \(B'\) 为 \(L \otimes_K E\) 的对应基，并设 \(U, V, W\) 为 \(u_{(L)}, v, w\) 关于 \(B'\) 的矩阵；注意， \(U\) 也是 \(u\) 关于 \(B\) 的矩阵，故其元素属于 \(K\) 。对每个 \(K\) -自同构 \(\sigma\) （即 \(L\) 的自同构），以及每个矩阵 \(A\) ，其元素属于 \(L\) ，令 \(A^\sigma\) 表示把 \(\sigma\) 作用于 \(A\) 的元素所得到的矩阵。设 \(\sigma\) 是一个 \(K\) -自同构，即 \(L\) 的自同构；则 \(U = U^\sigma = (V + W)^\sigma = V^\sigma + W^\sigma\) , \(V^\sigma W^\sigma = (VW)^\sigma = (WV)^\sigma = W^\sigma V^\sigma\) ；由于 \(V^\sigma\) 是某个可对角化自同态的矩阵，而 \(W^\sigma\) 是幂零的，由引理4可知 \(V^\sigma = V\) 且 \(W^\sigma = W\) 。由于这对所有 \(\sigma\) 都成立，V 和 W 的元素属于 K；若 \(u_{s}\) 与 \(u_{n}\) 是 E 的、关于基 B 的矩阵分别为 V 和 W 的自同态，则 \((u_{s})_{(L)} = v\) 且 \((u_{n})_{(L)} = w\) 。由此可知， \(u_{s}\) 是绝对半单的， \(u_{n}\) 是幂零的， \(u_{s}\) 与 \(u_{n}\) 可交换，并且 \(u = u_{s} + u_{n}\) 。这就完成了证明。

每当一个自同态 \(f\) 具有若尔当分解时，我们将其记为 \((f_s, f_n)\) ，并且自同态 \(f_s\) 与 \(f_n\) 分别称为 \(f\) 的绝对半单分量和幂零分量。当 \(K\) 是完全域时，每个自同态都有若尔当分解；在这种情况下，绝对半单自同态与半单自同态之间也没有区别，我们有时将“绝对半单分量”称为“半单分量”。

推论1. --- 设u有若尔当分解，且L是K的一个扩域。则 \(u_{(L)}\) 具有若尔当分解，其中 \((u_{(L)})_{s} = (u_{s})_{(L)}\) 且 \((u_{(L)})_{n} = (u_{n})_{(L)}\) .

推论2. --- 设u有若尔当分解。则E的每个与u可交换的自同态也与 \(u_{s}\) 和 \(u_{n}\) 可交换.

推论3. --- 设u和v是E的两个可交换的自同态，且它们具有若尔当分解。

\begin{enumerate}
\def\labelenumi{\alph{enumi})}
\item
  这些自同态 \(u, v, u_s, v_s, u_n, v_n\) 均交换。
\item
  这些自同态 \(u + v\) 与 \(uv\) 具有若尔当分解，且
\end{enumerate}

\[
\begin{array}{c} (u + v) _ {s} = u _ {s} + v _ {s}, \quad (u + v) _ {n} = u _ {n} + v _ {n}, \quad (u v) _ {s} = u _ {s} v _ {s}, \\ (u v) _ {n} = u _ {s} v _ {n} + u _ {n} v _ {s} + u _ {n} v _ {n}. \end{array}
\]

a) 部分由推论2得出。要证明 b) 部分，只需注意到 \(u_{s} + v_{s}\) 与 \(u_{s}v_{s}\) 是绝对半单的（VII，第43页，推论），且 \(u_{n} + v_{n}\) 与 \(u_{s}v_{n} + u_{n}v_{s} + u_{n}v_{n}\) 是幂零的（作为可交换的幂零自同态之和）。

推论4. --- 设u有若尔当分解，并设 R 是 K{[}X{]} 中的多项式。则自同态 \(\mathrm{R}(u)\) 有若尔当分解，且 \(\mathrm{R}(u)_{s} = \mathrm{R}(u_{s})\) .

注记. --- 1) 我们有 \(\det(u_s) = \det(u)\) 和 \(\operatorname{Tr}(u_s) = \operatorname{Tr}(u)\) .

\begin{enumerate}
\def\labelenumi{\arabic{enumi})}
\setcounter{enumi}{1}
\tightlist
\item
  u 可三角化的一个充分必要条件是：u 有一个若尔当分解，且 \(u_{s}\) 可对角化。则存在 E 的一组基，使得 u 关于该基的矩阵是下三角的，而 \(u_{s}\) 关于该基的矩阵是对角的，且与 u 的矩阵具有相同的对角线（参见下面的引理4和命题19）。
\end{enumerate}

但要注意，如果u关于某组基的矩阵是三角的，一般不能推出 \(u_{s}\) 关于同一组基的矩阵是对角的。

\begin{enumerate}
\def\labelenumi{\arabic{enumi})}
\setcounter{enumi}{2}
\tightlist
\item
  方阵的若尔当分解的概念可以类似地定义。例如，对于若尔当矩阵 \(U_{m,\alpha}\) 我们有
\end{enumerate}

\[
(U _ {m, \alpha}) _ {s} = \alpha . I _ {m}, (U _ {m, \alpha}) _ {n} = U _ {m, 0}.
\]

\begin{enumerate}
\def\labelenumi{\arabic{enumi})}
\setcounter{enumi}{3}
\tightlist
\item
  如果 \(u\) 是半单的但不是绝对半单的，则它没有若尔当分解。
\end{enumerate}

交换域上向量空间 V 的自同态 u 称为幂单的，如果自同态 \(u - Id_{V}\) 是幂零的，即存在整数 r 使得 \((u - Id_{V})^{r} = 0\) ；则 u 是 V 的自同构，因为若 \(u = Id_{V} - n\) ，其中 \(n^{r} = 0\) ，则

\[
\left(\operatorname{Id} _ {\mathrm{V}} + n + \dots + n ^ {r - 1}\right) u = u \left(\operatorname{Id} _ {\mathrm{V}} + n + \dots + n ^ {r - 1}\right) = \operatorname{Id} _ {\mathrm{V}}.
\]

若 V 的维数为 m，则 u 是幂单的当且仅当 \(\chi_{u}(X)=(X-1)^{m}\) （VII，第33页，命题5的推论3）。

命题17. --- 设E是交换域上的有限维向量空间，并设f是E的自同态。则以下条件等价：

\begin{enumerate}
\def\labelenumi{(\roman{enumi})}
\item
  \(f\) 有若尔当分解且是自同构；
\item
  \(f\) 有若尔当分解且 \(f_{s}\) 是自同构；
\item
  存在一个绝对半单的自同构 \(a\) 作用于 \(E\) ，和一个幂单的自同态 \(u\) 作用于 \(E\) ，使得 \(f = ua = au\) .
\end{enumerate}

此外，在这些条件下，在(iii)的记号中，我们必须有 \(a = f_{s}\) 和 \(u = 1 + f_{s}^{-1}f_{n}\) .

\begin{enumerate}
\def\labelenumi{(\roman{enumi})}
\item
  \(\Rightarrow\) （ii）：这由注记1得出。
\item
  \(\Rightarrow\) （iii）：取 \(a = f_{s}\) 和 \(u = 1 + f_{s}^{-1}f_{n}\) ；则 \(f = ua = au\) ，而 \(a\) 是一个绝对半单自同构，且 \(u\) 是幂单的。
\item
  \(\Rightarrow\) (i)：在(iii)的记号下，取 \(n = a(u - 1) = (u - 1)a\) ，则 \(an = na\) 且 \(f = a + n\) ，且 \(n\) 是幂零的。由此可知， \((a, n)\) 是 \(f\) 的若尔当分解。这蕴含了(i)以及关系式 \(a = f_s\) 和 \(u = 1 + f_s^{-1}f_n\) .
\end{enumerate}

令 \(f_{u}=f_{s}^{-1}f=ff_{s}^{-1}=1+f_{s}^{-1}f_{n}\) ，并称它为 f 的幂单分量。数对 \((f_{s},f_{u})\) 通常称为自同构 f 的乘法若尔当分解。

命题18. --- 设E是交换域K上的有限维向量空间，设u是E的一个自同态，并设E'是E中在u下封闭的子空间。设u'（分别地，u''）是由u诱导的E'（分别地，E/E'）的自同态。则 \(\chi_{u} = \chi_{u'} \cdot \chi_{u''}\) .

为使 u 具有若尔当分解，必要且充分条件是 \(u'\) 和 \(u''\) 具有；此外，若此成立，则 \(u'\) 和 \(u''\) 的绝对半单（分别地，幂零）分量是 \(E'\) 和 \(E/E'\) 的自同态，它们由 u 的绝对半单（分别地，幂零）分量诱导。

设 B 为 E 的一个基，它包含一个基 \(B'\) ，即 \(E'\) 的一个基，并设 \(B''\) 是

\(E'' = E/E'\) 的基，即 B - B' 的像。设 U, \(U'\) , \(U''\) 分别是 u, \(u'\) , \(u''\) 关于 B, \(B'\) , \(B''\) 的矩阵。则 U 具有形式

\[
\left( \begin{array}{c c} U ^ {\prime} & Z \\ 0 & U ^ {\prime \prime} \end{array} \right)
\]

且 \(\chi_{u} = \chi_{U} = \chi_{U'} \chi_{U''} = \chi_{u'} \chi_{u''}\) （参见 III，第 533 页，习题 2）。我们推断 u 的特征值集合是 \(u'\) 和 \(u''\) 的特征值集合的并。若 \(u'\) 和 \(u''\) 具有若尔当分解，则 \(u'\) 和 \(u''\) 的特征值在 K 上可分，因此 u 的特征值也是可分的，且 u 具有若尔当分解（VII，第43页，定理1）。反之，若 u 具有若尔当分解 (s, n)，则 s 和 n 保持 \(E'\) 不变，因为它们是 u 的多项式；设 \(s', n', s'', n''\) 表示 \(E', E', E'', E''\) 的、分别由 s、n、s、n 诱导的自同态。则 \(s'\) 和 \(s''\) 是绝对半单的，因为它们的极小多项式整除 s 的极小多项式（VII，第42页，命题15）；此外 \(n'\) 和 \(n''\) 是幂零的。最后 \(u' = s' + n', u'' = s'' + n'', s'n' = n's',\) 且 \(s''n'' = n''s''\) ，这就完成了证明。

命题19. --- 设E是交换域K上的有限维向量空间，且设F是由E的两两可交换的可三角化自同态组成的集合。则存在E的一组基，使得F中每个元素u关于这组基的矩阵都是下三角的，且 \(u_{s}\) 关于这组基的矩阵是对角的，其对角元素与u的相同。

根据VII第45页推论3，集合 \(\mathfrak{F}_s\) ——即 \(\mathfrak{F}\) 的元素的绝对半单分量组成的集合——由彼此交换的可对角化元素组成，故该集合可对角化（VII，第41页，命题13）；集合 \(\mathfrak{F}_n\) ——即 \(\mathfrak{F}\) 的元素的幂零分量组成的集合——由两两交换的幂零元素组成，且 \(\mathfrak{F}_n\) 的每个元素与 \(\mathfrak{F}_s\) 的每个元素交换。按照命题13（第七章，第41页）证明中的论证方式，我们看到存在将E分解为子空间的直和 \(E_i\) ，这些子空间在 \(\mathfrak{F}_s\) 和 \(\mathfrak{F}_n\) 下不变，并且使得 \(\mathfrak{F}_s\) 的每个元素在每一个 \(E_i\) 上的限制都是一个位似变换。将 E 依次替换为每个 \(E_i\) ，我们可以假设 \(\mathfrak{F}_s\) 的元素都是位似变换；只需证明存在 E 的一组基，使得 \(\mathfrak{F}_n\) 的元素在这组基下由对角线为零的下三角矩阵表示；这样我们就归结为 \(\mathfrak{F}\) 由幂零元素组成的情形。

现在假设 \(E \neq 0\) ，并设 F 为 E 的一个非零子空间，在 F 下不变，且维数最小。则对每个 \(u \in F\) ，u 在 F 上的限制的核非零且在 F 下不变（VII，第 41 页，引理 3）；由 F 的选取，u 在 F 上的限制从而对所有 \(u \in F\) 都为零。设 \(x \in F\) , \(x \neq 0\) ；则 \(u(x) = 0\) 对所有 \(u \in F\) 成立；通过对E的维数进行归纳论证，我们可以假设存在一个基 \((\bar{e}_{1}, \ldots, \bar{e}_{n-1})\) （即商 \(E' = E/Kx\) 的基），使得，对于所有 \(u \in F\) ，由 u 诱导的自同态 \(\bar{u}\) （ \(E'\) 上的）关于该基的矩阵是下三角的且对角线为零；如果 \(e_{i} \in E\) 投影到 \(\bar{e}_{i}\) 上，对于 i = 1， \ldots，n - 1，则基 \((e_{1}, \ldots, e_{n-1}, x)\) 满足所需条件。

\BourbakiExercises

\BourbakiExerciseGroup

\begin{enumerate}
\def\labelenumi{\arabic{enumi})}
\item
  证明：若 A 是整环，则多项式环 \(A[X_{\iota}]_{\iota \in I}\) 不是主理想整环，只要以下条件之一成立：1） \(\operatorname{Card}(I) \geqslant 2\) ；2） A 不是域（考虑理想 \((\mathbf{X}_{\alpha}) + (\mathbf{X}_{\beta})\) ，对于 \(\alpha \neq \beta\) ，以及 \((a) + (\mathbf{X}_{\alpha})\) ，其中 \(a \neq 0\) 是A的一个不可逆元素）。
\item
  证明在环 \(\mathbf{K}[[\mathbf{X}]]\) ——即域 \(\mathbf{K}\) 上的形式幂级数环——中，每个不可约元素都与 \(\mathbf{X}\) 相伴。
\item
  设 \(\mathbf{A}\) 是一个整环，其中每个非空理想集合都有极大元，并且每个极大理想都是主理想；证明 \(\mathbf{A}\) 是主理想整环。（注意，如果 \(\mathfrak{a} \neq \mathbf{A}\) 是 \(\mathbf{A}\) 的理想，则存在一个极大理想 \((p) \supset \mathfrak{a}\) 使得 \(\frac{1}{p} \mathfrak{a}\) ——它是 \(\mathbf{A}\) 的理想——严格包含 \(\mathfrak{a}\) .)
\end{enumerate}

\begin{enumerate}
\def\labelenumi{\alph{enumi})}
\setcounter{enumi}{1}
\tightlist
\item
  设 K 为一个域，设 \(\mathrm{K}_{1}=\mathrm{K}((\mathrm{X}))\) 为 K 上单未定元 X 的形式幂级数域（IV，第 38 页），并设 \(B=K_{1}[Y]\) 为系数在 \(K_{1}\) 中的、关于 Y 的形式幂级数环。设 A 为 B 的由那些幂级数 \(\sum_{n=0}^{\infty}a_{n}(\mathrm{X})\mathrm{Y}^{n}\) 组成的子环，其中幂级数 \(a_{0}(\mathrm{X})\) 仅包含指数 \(\geqslant 0\) 的 X 的幂。证明主理想 (X) 是 A 的唯一极大理想，但由元素 \(YX^{-n}\) ( \(n \geqslant 0\) 为任意整数）生成的 A 的理想不是主理想。
\end{enumerate}

\begin{enumerate}
\def\labelenumi{\arabic{enumi})}
\setcounter{enumi}{3}
\tightlist
\item
  设 A 是一个整环，并设 S 是 A 的一个不包含 0 的乘法封闭子集。设 \(S^{-1}A\) 表示由元素 \(s^{-1}a\) ( \(s \in S\) , \(a \in A\) ) 在 A 的分式域 K 中组成的环（I，第 116 页）。
\end{enumerate}

\begin{enumerate}
\def\labelenumi{\alph{enumi})}
\item
  证明如果 \(\mathbf{A}\) 是主理想整环，则 \(\mathbf{S}^{-1}\mathbf{A}\) 也是主理想整环.
\item
  设 A 为主理想整环，p 为 A 的不可约元素；证明理想 \((p)\) 的补集 S 是乘法封闭的；在环 \(S^{-1}A\) 中，理想 \(pS^{-1}A\) 是唯一的极大理想，且商域 \(S^{-1}A/pS^{-1}A\) 同构于 \(A/(p)\) .
\end{enumerate}

\begin{enumerate}
\def\labelenumi{\arabic{enumi})}
\setcounter{enumi}{4}
\tightlist
\item
  设 \(\mathbf{A}\) 是一个交换环，其中每个理想都是有限生成的。
\end{enumerate}

\begin{enumerate}
\def\labelenumi{\alph{enumi})}
\item
  元素 \(p \neq 0\) （\(\mathbf{A}\) 中的元素）称为不可除的，如果它不可逆，并且关系 \(p = xy (x, y \in \mathbf{A})\) 蕴含 \(x\) 或 \(y\) 属于 \(A p\) 。证明每个元素 \(a \neq 0\) （\(\mathbf{A}\) 中的元素）都可以（至少以一种方式）写成一些不可除元与一个可逆元素的乘积（利用第七章第4页的引理1）。
\item
  交换环A中的幂等元e被称为不可分解的，如果不存在一对非零幂等元f、g使得 \(f + g = e\) 且fg = 0。两个不同的不可分解幂等元的乘积为零。证明环A中的每个幂等元都是不可分解幂等元之和（证明若 \(e\) 是幂等元，则 \(1 - e\) 也是，并利用第七章第4页的引理1）。由此推出 \(\mathbf{A}\) 是有限个环 \(\mathbf{A}_i\) 的直积，其中每个环的每个理想都是有限生成的，且单位元是唯一的非零幂等元。假设这个族至少有两个元素。那么元素 \(a = \sum_{i} a_i\) （其中 \(a_i \in \mathbf{A}_i\) 对所有 \(i\) ) 是不可除的，当且仅当存在一个指标 \(k\) 使得
\end{enumerate}

该 \(a_{k}\) 在 \(\mathbf{A}_{k}\) 中是不可除的或为零，且 \(a_{i}\) 在 \(\mathbf{A}_{i}\) 中可逆（对所有 \(i \neq k\) ）；如果 \(a_{k} = 0\) ，则 \(\mathbf{A}_{k}\) 必然是一个整环。

\begin{enumerate}
\def\labelenumi{\alph{enumi})}
\setcounter{enumi}{2}
\item
  证明：若 1 是 \(\mathbf{A}\) 的唯一非零幂等元，且若 \(x \in \mathbf{A}\) 使得 \(x^k \in \mathbf{A}x^{k+1}\) 对于某个整数 \(k \geqslant 1\) 成立，则要么 \(x\) 可逆，要么 \(x^k = 0\) （若 \(x^k = ax^{k+1}\) ，则考虑 \(a^k x^k\) ).
\item
  设 \(p \neq 0\) 是 A 中的零因子。证明存在一个有限序列 \((a_k)_{1 \leqslant k \leqslant m}\) 由 A 的非零元素组成，使得 \(0 = pa_1\) , \(a_k = pa_{k+1}\) 对于 \(1 \leqslant k < m\) , \(a_m \notin \mathbf{A}_p\) ，以及 \((a_k) \neq (a_{k+1})\) 对于 \(1 \leqslant k < m\) .
\end{enumerate}

\begin{enumerate}
\def\labelenumi{\arabic{enumi})}
\setcounter{enumi}{5}
\tightlist
\item
  一个交换环 A 被称为主理想环，如果 A 的每个理想都是主理想。a) 证明主理想环的商环以及主理想环的有限乘积都是主理想环。
\end{enumerate}

\begin{enumerate}
\def\labelenumi{\alph{enumi})}
\setcounter{enumi}{1}
\item
  证明每个主理想环 A 都是某个有限族 \((\mathbf{A}_{i})_{1\leq i\leq m}\) 的直积，族中每个环都是主理想环，且其单位元是唯一的非零幂等元（利用习题5，b））。
\item
  设 A 是一个主理想环，其中单位元是唯一的非零幂等元，并且 A 含有除 0 以外的零因子。证明 A 中存在一个不可除的幂零元 p。（利用习题 5 a），证明存在一个不可除元 \(p \neq 0\) 属于 A，它是一个零因子；应用习题5 d)，然后证明主理想 Ab （\(y \in A\) 满足 \(yp^{m} \in Ap^{m+1}\) ）就是整个 A；最后应用习题5，c)。d) 在与c)相同的假设和记号下，假设 \(p^{m} = 0\) 且 \(p^{m-1} \neq 0\) 。证明 \(p^{m-1}\) 的零化子是 Ap（否则它将是一个主理想 Ac，满足 \(p \in Ac\) 且 \(c \notin Ap\) ；利用 p 不可除这一事实，推出 c 将是可逆的）；由此推出，对于 \(1 \leq k \leq m-1\) ， \(p^{m-k}\) 的零化子是 \(Ap^{k}\) 。由此得出理想 Ap 是极大的（仿照 d) 中第一个断言的论证，利用 1 - xp 在 A 中可逆，对所有 \(x \in A\) .
\item
  由 \(d\) ) 推出，对所有 \(x \in \mathbf{A}\) ，存在一个整数 \(k \geqslant 0\) 以及一个可逆元素 \(u\) 使得 \(x = up^k\) ，并且 \(\mathbf{A}\) 的理想只有诸 \(Ap^k\) ( \(0 \leqslant k \leqslant m\) ) （观察到 \(\mathbf{A}\) 中每个不属于 \(Ap\) 的元素都在 \(\mathbf{A}\) 中可逆).
\item
  设 A 为主理想环，令 \((x_{i})\) 为 A 的一族元素，并令 \(Ad\) 为由族 \((x_{i})\) 生成的理想；证明存在元素 \(x_{i}^{\prime}\) 属于 A，使得 \(x_{i} = dx_{i}^{\prime}\) 对所有 \(i\) 成立，以及 \(\sum_{i} Ax_{i}^{\prime} = A\) （化归到 1 是唯一的
\end{enumerate}

非零幂等元的情形）。

\(g)\) 给出一个主理想环和一个不可除元素 \(p\) （ \(\mathbf{A}\) 中的）的例子，使得理想 \(Ap\) 不是极大的（利用习题5， \(b\) )).

\begin{enumerate}
\def\labelenumi{\arabic{enumi})}
\setcounter{enumi}{6}
\tightlist
\item
  给定一个整环 \(\mathbf{A}\) ，我们称一个映射 \(w\) ——从 \(\mathbf{A}' = \mathbf{A} - \{0\}\) 到 \(\mathbf{N}\) 的映射——为欧几里得函数，如果它满足以下条件：
\end{enumerate}

\((S_{1}) w(xy) \geqslant w(y)\) 对于每一对 \(x, y\) ，其元素属于 \(\mathbf{A}'\) .

（S \(_{II}\) ) 如果 \(a \in \mathbf{A}'\) 且 \(b \in \mathbf{A}'\) ，则存在元素 \(q\) , \(r\) 属于 \(\mathbf{A}\) ，使得 \(a = bq + r\) ，并且要么 \(r = 0\) ，要么 \(w(r) < w(b)\) .

如果 A 上存在一个欧几里得函数，则称 A 为欧几里得整环。

\begin{enumerate}
\def\labelenumi{\alph{enumi})}
\item
  证明： \(\mathbf{Z}\) 和 \(\mathbf{K}[\mathbf{X}]\) （其中 \(\mathbf{K}\) 是域）都是欧几里得整环。
\item
  证明每个欧几里得整环都是主理想整环（如果 \(\mathfrak{a}\) 是 \(\mathbf{A}\) 的理想，选择 \(a \in \mathfrak{a}\) 使得 \(w(a)\) 尽可能小）。
\item
  证明高斯整数环 A（VII，第～1页）是关于函数 \(w(a + bi) = a^2 + b^2 = \mathbf{N}_{\mathbf{Q}(i) / \mathbf{Q}}(a + bi)\) 的欧几里得整环（注意，对所有 \(x \in \mathbf{Q}(i)\) ，存在 \(y \in \mathbf{A}\) 使得 \(\mathbf{N}(x - y) \leqslant 1/2\) ).
\end{enumerate}

\(d\) ）证明二次扩张环 \(\mathbf{B}\) （\(\mathbf{Z}\) 的、以 \((1,e)(e^2 = 2)\) 为基的环）是关于函数

\[
w (a + b e) = \left| a ^ {2} - 2 b ^ {2} \right| = \left| N _ {Q (e) / Q} (a + b e) \right|
\]

（对所有 \(x \in \mathbf{Q}(e)\) ，存在 \(y \in B\) 使得 \(|\mathbf{N}(x - y)| \leqslant 1/2\) )的欧几里得整环.

\begin{enumerate}
\def\labelenumi{\arabic{enumi})}
\setcounter{enumi}{7}
\tightlist
\item
  设 \(\mathbf{K}\) 是一个二次域，即 \(\mathbf{Q}\) 的次数为 2 的扩张。
\end{enumerate}

\begin{enumerate}
\def\labelenumi{\alph{enumi})}
\item
  证明： \(\mathbf{K} = \mathbf{Q}(e)\) ，其中 \(e^2\) 是一个有理整数 \(d \neq 1\) ，不被 \(>1\) 的平方数整除（在 \(\mathbf{Z}\) 中）.
\item
  如果 \(\alpha = a + be\) ( \(a, b \in \mathbf{Q}\) ），令 \(\overline{\alpha}\) 表示共轭 \(a - be\) （即 \(\alpha\) 的共轭）。我们称 \(\alpha\) 是二次域 \(\mathbf{K}\) 的整数，如果其范数 \(\alpha\overline{\alpha}\) 及其迹 \(\alpha + \overline{\alpha}\) 是有理整数。证明这些整数（\(\mathbf{K}\) 的整数）构成一个环 \(\mathbf{A}\) （注意如果 \(\alpha\) 和 \(\beta\) 是 \(\mathbf{K}\) 的整数，那么有理数 \(\alpha\beta + \overline{\alpha\beta}\) 和 \(\alpha\overline{\beta} + \overline{\alpha}\beta\) 的和与积是整数，因此这些有理数是 \(\mathbf{Z}\) 上某个首一多项式的根；由此推断它们是有理整数）。
\item
  证明 \(\mathbf{K} = \mathbf{Q}(e)\) 的整数构成一个 \(\mathbf{Z}\) -模，其基为 \((1, e)\) ，当 \(d - 1 \not\equiv 0 \pmod{4}\) 时；或为 \((1/2(1 + e), 1/2(1 - e))\) ，当 \(d - 1 \equiv 0 \pmod{4}\) 时（若 \(a + be\) 是 \(\mathbf{K}\) 的整数，验证 \(2a\) 和 \(2b\) 是满足 \((2a)^2 - d(2b)^2 \equiv 0(4)\) 的有理整数，并研究它们是否可能为奇数）。证明这些基的判别式（III，第～549页，定义～3）分别为 \(4d\) 和 \(d\) 。
\end{enumerate}

\begin{enumerate}
\def\labelenumi{\arabic{enumi})}
\setcounter{enumi}{8}
\item
  *（闵可夫斯基定理）若 \(a, b, c, d\) 是实数，且 \(\mathbf{D} = \begin{vmatrix} a & b \\ c & d \end{vmatrix}\) ，则不等式 \(|na + mb| \leqslant \mathbf{A}\) , \(|nc + md| \leqslant \mathbf{B}\) 有非零整数解 \((n, m)\) ，只要 \(\mathbf{A} > 0\) , \(\mathbf{B} > 0\) , \(\mathbf{D} > 0\) 且 \(\mathbf{AB} \geqslant \mathbf{D}\) ，且这样的解只有有限多个（考虑离散子群 \(\mathbf{G}\) ——\(\mathbf{R}^2\) 中由 \((a, c)\) 和 \((b, d)\) 生成的子群——，并证明矩形 \(\mathbf{E}\) ——由关系 \(0 \leqslant x \leqslant \mathbf{A}\) , \(0 \leqslant y \leqslant \mathbf{B}\) 定义，其面积大于在向量 \((a, c)\) 和 \((b, d)\) 上构造的平行四边形的面积——包含两个模 \(\mathbf{G}\) 同余的不同向量；用反证法论证，注意到否则 \(p^2\) 个矩形——它们由 \(\mathbf{E}\) 借助平移（ \(ha + kb\) , \(hc + kd\) ）（其中 \(0 \leqslant h < p\) 且 \(0 \leqslant k < p\) ）得到——将是两两不相交的）。
\item
  设 \(d\) 为一个不被 \(>1\) 的平方数整除的有理整数。在二次域 \(\mathbf{K} = \mathbf{Q}(\sqrt{d})\) 中，我们将考察单位，即环 \(\mathbf{A}\) ——\(\mathbf{K}\) 的整数环——中的可逆元素（习题8，b））。
\end{enumerate}

\begin{enumerate}
\def\labelenumi{\alph{enumi})}
\item
  如果 \(d < 0\) （“虚二次域”），则证明 \(K\) 的单位只有 1 和 -1，但 \(d = -1\) 例外，此时 \(i\) 和 -i 也是单位；以及 \(d = -3\) 的情形，此时四个数 \(1/2(\pm 1 \pm \sqrt{-3})\) 也是单位。
\item
  从现在起假设 \(d > 0\) （“实二次域”），并假设 \(\mathbf{K}\) 嵌入到 \(\mathbf{Q}\) 的一个极大序扩张中 (VI, 第～25页)，* 例如在 \(\mathbf{R}\) 中. * 证明正单位构成 K 的单位群 U 的子群 \(U'\) ，并且 U 是 \(U'\) 与 \(\{-1,1\}\) 的直积.
\item
  证明 K 的单位是范数为 1 或 -1 的整数。
\end{enumerate}

\(d\) ) 证明群 \(\mathbf{U}^{\prime}\) 是非平凡的（令 D 表示整数基的判别式（习题 8， \(c\) )；证明 K 中的整数 \(\alpha\) （满足 \(\mathbf{N}(\alpha)\leqslant \mathbf{D}\) ）被划分为模 U 的有限多个同余类；在每个类中选取一个整数 \(\beta_{i}\) ；如果 B 是一个有理数，使得 \(0 < \mathrm{B} < \inf (|\beta_i|)\) ，证明（习题9）存在 K 的一个整数 \(\beta\) ，满足 \(0 < |\beta |\leqslant \mathrm{B}\) 且 \(\mathbf{N}(\beta)\leqslant \mathbf{D}\) ；由此推出某个 \(\beta \beta_{i}^{-1}\) 是一个单位 \(\eta\) ，满足 \(|\eta| < 1\) ).

\begin{enumerate}
\def\labelenumi{\alph{enumi})}
\setcounter{enumi}{4}
\item
  证明： \(\mathbf{U}'\) 同构于 \(\mathbf{Z}\) （证明存在一个单位 \(\varepsilon\) ，使得 \(|\varepsilon| < 1\) 且 \(|\varepsilon|\) 在具有此性质的数中尽可能大，利用习题9以及事实 \(|\overline{\xi}| = |\xi|^{-1}\) （对单位成立）。
\item
  证明在 \(\mathbf{Q}(\sqrt{2})\) 中，元素 \(\sqrt{2} - 1\) 是 \(\mathbf{U}'\) 的一个生成元，且其范数为 \(-1\) .
\end{enumerate}

\(g\) ) 证明如果 \(d\) 有一个素因子 \(p\) （形如 \(4n - 1\) ），则 \(\mathbf{U}'\) 的每个生成元都具有范数 1。（注意， \(-1\) 在域 \(\mathbf{F}_p\) 中不能是平方数，利用 V，第 93 页，命题 1，c)。

¶ 11) 我们研究这样的二次域 \(\mathbf{Q}(\sqrt{d})\) ：映射 \(\alpha \mapsto |\mathrm{N}(\alpha)|\) 是该域整数上的欧几里得函数（习题 7）。

\begin{enumerate}
\def\labelenumi{\alph{enumi})}
\item
  证明其充分必要条件为：对所有 \(\alpha \in \mathbf{Q}(\sqrt{d})\) ，存在该域的一个整数 \(\beta\) 使得 \(|\mathrm{N}(\alpha - \beta)| < 1\) ；写出此条件，将 \(\alpha\) 和 \(\beta\) 用该域的整数基表示（习题8，c））。
\item
  对于虚二次域的情形 \((d < 0)\) ，证明使 \(|\mathbf{N}(\alpha)|\) 是欧几里得函数的 d 的值只有 -1、-2、-3、-7、-11。
\item
  对于实二次域，我们将仅限于 \(d \equiv 2\) 或 3（模 4）的情形。证明此时仅有有限多个 \(d\) 的值使得 \(|\mathrm{N}(\alpha)|\) 是欧几里得函数。（若 \(|\mathrm{N}(\alpha)|\) 是欧几里得函数，且 \(a\) 是一个自然数，使得 \(0 \leqslant a \leqslant d\) 且 \(a \equiv t^2(d)\) ，则证明 \(a\) 或 \(d - a\) 之一具有形式 \(x^2 - dy^2\) ，其中 \(x\) 和 \(y\) 是有理整数，这通过考虑元素 \(t / \sqrt{d}\) （属于 \(\mathbf{Q}(\sqrt{d})\) ）来证明；接下来证明，对于足够大的 \(d\) ，存在一个奇数 \(z\) 使得 \(5d < z^2 < 6d\) （若 \(d \equiv 3 \mod 4\) ）或 \(2d < z^2 < 3d\) （若 \(d \equiv 2 \mod 4\) ）；然后取 \(a = z^2 - d[z^2/d]\) ，并证明在这些条件下，\(a\) 与 \(d - a\) 都不具有形式 \(x^2 - dy^2\) （通过考察模 8 的剩余类）。证明 \(d\) 的值，当这样的整数 \(z\) 不存在时，只有 3、7、11、19、35、47 和 59（对于 \(d \equiv 3 \mod 4\) ）以及 2、6、14 和 26（对于 \(d \equiv 2 \mod 4\) ）。证明对于 \(d = 2\) 和 \(d = 3\) 确实存在欧几里得函数。
\end{enumerate}

\begin{enumerate}
\def\labelenumi{\arabic{enumi})}
\setcounter{enumi}{11}
\tightlist
\item
  证明二次域 \(\mathbf{Q}(\alpha)\) （其中 \(\alpha^2 = -5\) ）的整数环不是主理想整环（参见第六章，第34页，习题27）。对 \(\mathbf{Q}(\beta)\) （其中 \(\beta^2 = 10\) ）提出同样的问题（我们有
\end{enumerate}

\[
2. 3 = (4 + \beta) (4 - \beta)).
\]

\begin{enumerate}
\def\labelenumi{\arabic{enumi})}
\setcounter{enumi}{12}
\item
  设 \(f\) 是 \(\mathbf{Z}\) 上一个未定元的非常数多项式。证明对于每个整数 \(k \geqslant 1\) ，存在无穷多个整数 \(n \in \mathbf{N}\) 使得 \(f(n)\) 非零且至少被 \(k\) 个不同的素数整除（可以对 \(k\) 用归纳法：如果 \(n \in \mathbf{N}\) 使得 \(f(n)\) 非零且至少被 \(k\) 个不同的素数整除，则考虑多项式 \(f(n + Xf(n)^2)\) .
\item
  形如 \(4n - 1\) （相应地 \(6n - 1\) ）的素数集合是无限的（方法与命题 5（VII，第～5 页）相同；注意除 2（相应地，除 2 和 3）以外的每个素数都具有形式 \(4n \pm 1\) （相应地 \(6n \pm 1\) )).
\item
  形如 \(2^k + 1\) ( \(k \geqslant 1\) 为整数）的数只有当 \(k\) 是 2 的幂时才可能是素数。如果素数 \(p\) 整除 \(2^{2^n} + 1\) ，那么它不能整除 \(2^{2^m} + 1\) ，对于 \(m > n\) ( \(2^{2^m} \equiv 1(p)\) ）；利用这一点给出命题5的一个新证明。证明 \(2^{2^5} + 1\) 不是素数（令 \(a = 2^7\) 和 \(b = 5\) ；证明 \(1 + ab - b^4 = 2^4\) 且 \(1 + ab\) 整除
\end{enumerate}

\[
n = 2 ^ {4} a ^ {4} + 1 = (1 + a b - b ^ {4}) a ^ {4} + 1.
\]

\begin{enumerate}
\def\labelenumi{\arabic{enumi})}
\setcounter{enumi}{15}
\item
  形如 \(a^k - 1\) ( \(a\) 和 \(k\) 为整数， \(a > 0\) , \(k > 1\) ) 的数只有当 \(a = 2\) 且 \(k\) 是素数时才可能是素数。
\item
  证明如果 \(a^n + 1\) 是素数，其中 \(a > 1\) 且 \(n > 1\) ，则 \(a\) 是偶数且 \(n = 2^m\) .
\item
  设 M 是由有限个整数 \(q_{i} > 1\) ( \(1 \leqslant i \leqslant m\) ) 生成的乘法幺半群。证明：对于每个整数 k \textgreater{} 0，存在 k 个连续整数，它们都不属于 M（如果 \(h_{1}, \ldots, h_{m}\) 是 m 个 \textgreater{} 1 的整数，且 k 是都不被任何 \(h_{i}\) 整除的连续整数的最大个数，则 \((k + 1)^{-1} \leqslant h_{1}^{-1} + \cdots + h_{m}^{-1}\) ；现在注意到，对所有 r \textgreater{} 0，M 中每个足够大的数都至少被一个 \(q_{i}^{+}\) 整除）。由此结果给出 VII 第5页命题5的一个新证明。
\item
  对每个整数 \(n > 0\) ，证明素数 \(p\) 在 \(n!\) 的素因子分解中的指数是 \(\sum_{k=1}^{\infty} [np^{-k}]\) （其中只有有限多个非零项）。
\item
  * 设 \(\pi(x)\) 为小于实数 \(x > 0\) 的素数的个数。对于每个整数 \(n > 0\) ，证明
\end{enumerate}

\[
n ^ {\pi (2 n) - \pi (n)} \leqslant \binom {2 n} {n} \leqslant (2 n) ^ {\pi (2 n)}
\]

（注意 \(\binom{2n}{n}\) 是所有素数 \(q\) （满足 \(n < q \leqslant 2n\) ）的乘积的倍数）；另一方面，利用习题19证明 \(\binom{2n}{n}\) 整除乘积 \(\prod p^{r(p)}\) ，其中 \(p\) 遍历素数集合 \(\leqslant 2n\) ，而 \(r(p)\) 是满足 \(p^{r(p)} \leqslant 2n\) 的最大整数。由此结果推出存在两个实数 \(a\) 和 \(b\) 使得 \(0 < a < b\) 且

\[
a n (\log n) ^ {- 1} \leqslant \pi (n) \leqslant b n (\log n) ^ {- 1}
\]

（注意 \(2^{2n}(2n + 1)^{-1} \leqslant (2n)! (n!)^{-2} \leqslant 2^{2n}\) ). *

\(^{1}\) 已证明 \(\pi(x) \sim \frac{x}{\log x}\) 当 x 趋于 \(+\infty\) 时（参见例如 A. E. INGHAM，《素数的分布》（剑桥丛刊，第 30 号），剑桥大学出版社，1932 年）。

\begin{enumerate}
\def\labelenumi{\arabic{enumi})}
\setcounter{enumi}{20}
\tightlist
\item
  证明对于 \(m \geqslant 1\) 和 \(n \geqslant 1\) ，有理数 \(\sum_{k=n}^{n+m} k^{-1}\) 绝不是整数（若
\end{enumerate}

\(2^{q}\) 是整除数 \(n, n+1, \ldots, n+m\) 中至少一个的 2 的最大幂，则它只整除这些数中的一个）。

\begin{enumerate}
\def\labelenumi{\arabic{enumi})}
\setcounter{enumi}{21}
\item
  设 \(a\) 和 \(b\) 是两个 \(> 0\) 的整数，设 \(d\) 是 \(a\) 和 \(b\) 的最大公因数，并设 \(n\) 是一个 \(\geqslant 1\) 的整数；令 \(a = da_1\) 和 \(b = db_1\) ；证明分数 \(a(a + b) \ldots (a + (n - 1)b) / n!\) 的既约表达式中的分母只具有整除 \(b_1\) 且 \(\leqslant n\) 的素因子（注意，如果 \(p\) 是一个不整除 \(b_1\) 的素数，则 \(b_1\) 模 \(p^r\) 可逆，对所有 \(r > 0\) ). 由此直接证明二项式系数 \(n! / (m!(n - m)! )\) 是整数。
\item
  设 \(n > 1\) 是一个有理整数；证明这 \(n - 1\) 个二项式系数 \(\binom{n}{k}\) ( \(1 \leqslant k \leqslant n - 1\) ) 的最大公因数等于 1，如果 \(n\) 被两个不同的素数整除；等于 \(p\) ，如果 \(n\) 是素数 \(p\) 的幂。
\item
  \begin{enumerate}
  \def\labelenumii{\alph{enumii})}
  \tightlist
  \item
    设 \(n = \prod_{k} p_k^{\nu_k} > 0\) 是一个用其素因子表示的有理整数。
  \end{enumerate}
\end{enumerate}

证明 n 的正因数之和由公式给出

\[
\sigma (n) = \prod_ {k} \frac {p _ {k} ^ {\nu_ {k} + 1} - 1}{p _ {k} - 1}.
\]

\begin{enumerate}
\def\labelenumi{\alph{enumi})}
\setcounter{enumi}{1}
\tightlist
\item
  我们说 \(n\) 是完全数，如果 \(2n = \sigma(n)\) 。证明一个偶数 \(n\) 是完全数，当且仅当它具有形式 \(2^h(2^{h+1} - 1)\) ，其中 \(h \geqslant 1\) 且 \(2^{h+1} - 1\) 是素数。（如果 \(n\) 具有形式 \(2^m q\) ，其中 \(m \geqslant 1\) 且 \(q\) 是奇数，证明 \(q\) 能被 \(2^{m+1} - 1\) 整除，并由此推出，如果 \(q\) 不是素数，那么我们将有 \(\sigma(n) > 2n\)）。
\end{enumerate}

\begin{enumerate}
\def\labelenumi{\arabic{enumi})}
\setcounter{enumi}{24}
\item
  设 \(a\) 和 \(b\) 是两个 \(> 0\) 的互素整数。若 \(a\) 和 \(b\) 都 \(> 2\) ，则存在整数 \(x\) 和 \(y\) 使得 \(ax + by = 1\) 且 \(|x| < 1/2b\) , \(|y| < 1/2a\) ；此外，整数 \(x\) 和 \(y\) 由这些条件唯一确定。当 \(a, b\) 之一等于 2 时情况如何？
\item
  设 \(x, y\) 是两个互素的整数；则每个整数 \(z\) （满足 \(|z| < |xy|\) ）都可以以一种或两种方式表示为 \(z = ux + vy\) 的形式，其中 \(u\) 和 \(v\) 是满足 \(|u| < |y|\) 和 \(|v| < |x|\) 的整数。若 \(z\) 能以这种形式用两种不同方式表示，则 \(x\) 和 \(y\) 都不能整除 \(z\) ；例子 \(x = 5, y = 7, z = 12\) 表明其逆命题不成立。
\item
  设 \(f\) 和 \(g\) 是 \(\mathbf{K}[\mathbf{X}]\) 中的两个互素多项式；证明每个多项式 \(h\) （满足 \(\deg(h) < \deg(fg)\) ）都可以唯一地表示为 \(h = uf + vg\) 的形式，其中 \(u\) 和 \(v\) 是满足 \(\deg(u) < \deg(g)\) 和 \(\deg(v) < \deg(f)\) 的多项式。首先给出一个与习题26类似的证明（基于欧几里得算法），然后给出一个将该问题视为线性问题的证明。
\item
  \begin{enumerate}
  \def\labelenumii{\alph{enumii})}
  \tightlist
  \item
    设K是一个代数闭的交换域，设f是K{[}X{]}中的一个非常数的首一多项式，并设g是K{[}X, Y{]}中的一个多项式；假设对于f的每一个根 \(\alpha_{i}\) ，都存在一个根 \(\beta_{i}\) ——即 \(g(\alpha_{i}, \mathbf{Y}) = 0\) 的根——使得 \(\frac{\partial g}{\partial \mathbf{Y}} (\alpha_{i}, \beta_{i}) \neq 0\) 。则证明存在一个多项式 \(h(\mathbf{X}) \in \mathbf{K}[\mathbf{X}]\) 使得 \(g(\mathbf{X}, h(\mathbf{X})) \equiv 0 \pmod{f(\mathbf{X})}\) （利用命题4，化归到情形 \(f(\mathbf{X}) = \mathbf{X}^{m}\) （对某个 \(m \geqslant 1\) ）；注意环 \(K[X]/(X^{m})\) 同构于商环 \(K[[X]]/(X^{m})\) ，并利用 IV 第 37 页的推论）。
  \end{enumerate}
\end{enumerate}

\begin{enumerate}
\def\labelenumi{\alph{enumi})}
\setcounter{enumi}{1}
\tightlist
\item
  特别地证明，若 m 是不被 K 的特征整除的整数，则对每一对互素的首一多项式 f 和 \(f_{0}\) （在 K{[}X{]} 中），存在一个多项式 \(h(\mathbf{X})\) （在 K{[}X{]} 中），使得 \((h(\mathbf{X}))^{m} \equiv f_{0}(\mathbf{X}) \pmod{f(\mathbf{X})}\) 。
\end{enumerate}

\BourbakiExerciseGroup

\begin{enumerate}
\def\labelenumi{\arabic{enumi})}
\item
  设 A 为主理想整环，设 M 为挠 A-模，设 x, y 为 M 的两个元素，并设 Aα 与 Aβ 分别为 x, y 的零化子；若 δ 为 α 与 β 的一个最大公因子，证明 \(x + y\) 的零化子 Aγ 使得 γ 整除 α 与 β 的最小公倍数 αβ/δ，并且是 αβ/δ² 的倍数。当 α 与 β 是 A 的同一个不可约元素 π 的不同幂 \(\pi^{\mu}\) , \(\pi^{\nu}\) 时，证明 \(\gamma = \pi^{\sup(\mu, \nu)}\) 。
\item
  设 A 为主理想整环，设 \(\pi\) 为 A 的一个不可约元素，并设 M 为 \(\pi\) -准素模；证明对所有 \(x \in M\) 以及所有 \(\alpha \in A\) （不是 \(\pi\) 的倍数），存在唯一的 \(y \in M\) 使得 \(x = \alpha y\) （利用贝祖恒等式）；令 \(y = \alpha^{-1}x\) 。由此推出，如果 \(A_{(\pi)}\) 表示由 A 的分式域中形如 \(\beta/\alpha\) 、满足 \(\alpha, \beta \in A\) 且 \(\alpha\) 不是 \(\pi\) 的倍数的元素构成的主理想整环（VII，第48页，习题4），则 M 上存在唯一的 \(A_{(\pi)}\) -模结构，使得算子环到 A 的限制给出 M 上原有的 A-模结构。这样定义的 \(A_{(\pi)}\) -模被称为与 M 典范关联；其子模与 M 的 A-子模完全相同。
\item
  设 A 为主理想整环。
\end{enumerate}

\begin{enumerate}
\def\labelenumi{\alph{enumi})}
\item
  一个 A-模 M 是内射的（II，第389页，习题11），当且仅当对于所有 \(x \in \mathbf{M}\) 以及 A 中所有 \(\alpha \neq 0\) ，存在 \(y \in \mathbf{M}\) 使得 \(x = \alpha y\) ；此时称 A-模 M 是可除的。若 E 是任意 A-模，则 E 的所有可除子模之和是 E 的最大可除子模 \(\Delta(E)\) 。对 E 的每个子模 F，我们有 \(\Delta(F) \subset \Delta(E)\) 。
\item
  设 M 是一个可除 A-模，设 \(\pi\) 为 A 的一个不可约元素，且令 \(x_{0}\) 是 M 中的一个元素，其零化子具有形式 \(A\pi^{n}\) (n \textgreater{} 0 为整数）。递归地定义 M 的元素的一个序列 \((x_{s})\) ，条件为 \(x_{s} = \pi x_{s+1}\) 对每个整数 \(s \geqslant 0\) 成立。证明由序列 \((x_{s})\) 生成的 M 的子模同构于模 U = K/A 的 \(\pi\) -准素分量 \(U_{\pi}\) ，其中 K 是 A 的分式域。c) 证明 \(U_{\pi}\) 是一个可除模，并且 \(U_{\pi}\) 的每个不同于 \(U_{\pi}\) 和 0 的子模都是循环的，且由 \(\pi\) 的某个幂关于 mod A 的类生成；由此推出 \(U_{\pi}\) 是不可分解的（II，第393页，习题21）。
\end{enumerate}

\(d\) ）证明每个不可分解的可除 \(\mathbf{A}\) -模同构于 \(\mathbf{K}\) ，或同构于诸模 \(\mathrm{U}_{\pi}\) 之一（利用 \(a\) ）和 \(b)\) ）。

\begin{enumerate}
\def\labelenumi{\alph{enumi})}
\setcounter{enumi}{4}
\tightlist
\item
  证明每个可除A-模都同构于一族不可分解可除子模的直和（将Zorn引理应用于其和为直和的不可分解可除子模族）。
\end{enumerate}

\begin{enumerate}
\def\labelenumi{\arabic{enumi})}
\setcounter{enumi}{3}
\tightlist
\item
  一个族 \((x_{i})\) （由模 \(\mathbf{M}\) ——主理想整环 \(\mathbf{A}\) 上的——的非零元素组成）称为拟自由的，如果每个形如 \(\sum_{i} \alpha_{i} x_{i} = \beta y\) 的关系式 ( \(y \in \mathbf{M}, \alpha_{i} \in \mathbf{A}\) ,
\end{enumerate}

\(\beta \in \mathbf{A}\) ) 都蕴含存在一族元素 \(\alpha_{i}^{\prime} \in \mathbf{A}\) ，使得 \(\alpha_{i}x_{i} = \beta \alpha_{i}^{\prime} x_{i}\) 对每个指标 \(i\) 成立。于是关系式 \(i \neq k\) 蕴含 \(x_{i} \neq x_{k}\) ；拟自由族的元素组成的集合称为 \(\mathbf{M}\) 的拟自由子集。

\begin{enumerate}
\def\labelenumi{\alph{enumi})}
\item
  设 \(\pi\) 是 A 的一个不可约元素。证明：在具有零化子 \(\mathbf{A}\pi^n\) ( \(n \geqslant 1\) ) 的 A-模中，任何具有零化子 \(\mathbf{A}\pi^n\) 的元素都是拟自由的（使用贝祖恒等式）。
\item
  设 \((x_{i})\) 是 A-模 M 中的一个拟自由族，并设 N 为由该族生成的 M 的子模。证明对 M/N 的每个元素 \(\dot{x}\) ，存在一个元素 \(x \in \dot{x}\) ，其零化子等于 \(\dot{x}\) 在 M/N 中的零化子。若 \(\dot{x}\) 在 M/N 中还是拟自由的，证明由 \((x_{i})\) 以及 x 组成的族在 M 中是拟自由的。
\item
  由 a) 和 b) 可推出，若 A-模 M 的零化子非零，则 M 是循环模的直和（化归到 \(\pi\) -准素模的情形，并对 M 的拟自由元素族应用 Zorn 引理）（参见 VII，第 19 页，定理 2）。
\end{enumerate}

\begin{enumerate}
\def\labelenumi{\arabic{enumi})}
\setcounter{enumi}{4}
\item
  设 M 是主理想整环 A 上的有限生成挠模。证明 M 是有限长度的模；因此，M 的任意非空子模集合都有极大元和极小元。
\item
  设 M 是主理想整环 A 上的挠模，它不是有限生成的，但其所有真子模都是有限生成的。证明存在 A 的一个不可约元素 \(\pi\) ，使得 M 同构于模 U = K/A 的 \(\pi\) -准素分量 \(U_{\pi}\) （习题 3，c））。（首先证明 M 的非零 \(\pi\) -准素分量不能多于一个，因而 M 是一个 \(\pi\) -准素模，对某个不可约元素 \(\pi \in A\) 。然后证明 M = \(\pi M\) ，这可通过观察 M/ \(\pi M\) 若非零则不可能有限生成得到；但 M/ \(\pi M\) 是域 A/ \(\pi A\) 上的向量空间，并且会包含一个非有限生成的真子空间。）
\item
  设 M 是主理想整环 A 上的模；M 的子模 N 被称为纯的，如果对所有 \(\alpha \in A\) ，则对 \(N \cap (\alpha M) = \alpha N\) .
\end{enumerate}

\begin{enumerate}
\def\labelenumi{\alph{enumi})}
\item
  M 的子模 N 是纯的，当且仅当对于所有 \(\alpha \in A\) 和每个元素 \(\dot{x} \in M/N\) （其零化子为 \(A\alpha\) ），存在一个元素 \(x \in \dot{x}\) ，其零化子为 \(A\alpha\) 。
\item
  证明 M 的挠子模是纯的。 M 的元素的族 \((x_{t})\) 是拟自由的（习题 4），当且仅当子模 \(Ax_{t}\) 的和是直和并且是 M 的纯子模。证明如果 P 和 Q 是 M 的子模，使得 \(P \cap Q\) 和 \(P + Q\) 是纯的，则 P 和 Q 是纯的。如果 Q 是 M 的纯子模，且 \(P \supset Q\) 是 M 的子模，则 P 在 M 中是纯的，当且仅当 P/Q 在 M/Q 中是纯的。
\item
  证明 M 的每个作为 M 的直因子的子模 N 都是纯的。对于 M 的直因子的滤过族的并集，同样成立。反之，如果 N 是 M 的纯子模，且 M/N 是循环子模的直和，证明 N 是 M 的直因子（使用 a））。给出 M 的直因子 P 和 Q 的例子，使得 \(P \cap Q\) （相应地 \(P + Q\) ）不是纯子模（对于第一个例子，取 \(M = (\mathbf{Z}/4\mathbf{Z}) \oplus (\mathbf{Z}/2\mathbf{Z})\) ，对于第二个例子，取 \(M = Z^{2}\) ）。
\item
  设 N 是 M 的纯子模，其零化子 Aα 非零。如果 f 是从 M 到 M/αM 的自然同态，证明 f 限制为从 N 到 f(N) 的同构，并且 f(N) 是 M/αM 的纯子模（使用 Bézout 恒等式）。由此推出（借助 c）和习题 4，c）） \(f(N)\) 是 M/ \(\alpha M\) 中的直因子，并且 N 是 M 中的直因子。特别地，如果模 M 的挠子模具有非零零化子（当它是有限生成时就是这种情况），则它是 M 中的直因子（参见 VII，第 20 页，定理 2 的推论 1）。
\item
  设 \(p\) 是一个素数；设 \(\mathbf{N}\) 表示循环 \(\mathbf{Z}\) -模 \(\mathbf{Z} / (p^n)\) ( \(n \geqslant 1\) ) 的直和，并令 \(e_n\) 表示 \(\mathbf{N}\) 中这样的元素：其所有坐标均为零，除了第 \(n\) 个坐标，它等于 1 模 \(p^n\) 的类。设 \(\mathbf{M}\) 为 \(\mathbf{Z}\) -模 \(\mathbf{N} \times \mathbf{Q}\) 的由元素 \((e_n, p^{-n})\) 生成的子模。证明挠子模 \(\mathbf{P}\) （ \(\mathbf{M}\) 的）不是 \(\mathbf{M}\) 中的直因子（注意， \(\mathbf{M} / \mathbf{P}\) 的每个元素都能被 \(p^n\) 整除，对所有 \(n\) ，并且 \(\mathbf{M}\) 中没有元素具有此性质，除了 \((0, 1)\) 的倍数）。
\end{enumerate}

\begin{enumerate}
\def\labelenumi{\arabic{enumi})}
\setcounter{enumi}{7}
\tightlist
\item
  设 A 为主理想整环，设 \(\pi\) 为 A 的一个不可约元素，并设 M 为 \(\pi\) -准素模。则 M 的元素 \(x\) 称为具有高度 \(n\) ，如果 \(x \in \pi^n\mathbf{M}\) 且 \(x \notin \pi^{n+1}\mathbf{M}\) 。若 \(x \in \pi^n\mathbf{M}\) 对所有 \(n \geqslant 1\) 成立，则称 \(x\) 具有无限高度。
\end{enumerate}

\begin{enumerate}
\def\labelenumi{\alph{enumi})}
\item
  证明M是可除的当且仅当其所有元素都具有无限高度（习题3）。
\item
  设 x 为 M 中高度为 n 的元素，使得 \(\pi x = 0\) ；证明若 \(x = \pi^{n}y\) 则 M 的循环子模 Ay 是 M 的一个直和因子（注意 Ay 是纯子模，并应用习题 7，d））。
\item
  证明一个不可除的 \(\pi\) -准素模 M 总是包含至少一个非零循环直和因子（若 \(x \in \mathbf{M}\) 具有有限高度，且存在整数 \(m\) 使得 \(\pi^m x \neq 0\) 具有无限高度，则令 \(s\) 为满足该性质的最小整数；写出 \(\pi^s x = \pi y\) ，其中 \(y\) 的高度严格大于 \(\pi^{s-1}x\) ，并将 \(b\) ) 应用于 \(y - \pi^{s-1}x\) ）。由此推出，任何不可分解的 \(\pi\) -准素模要么是循环的，要么同构于一个模 \(\mathrm{U}_{\pi}\) （习题3）。d) 证明：若 \(\pi\) -准素模 M 的每个纯子模都是 M 中的直和因子，则 M 是一个可除模 P 与一个具有非零零化子的 \(\pi\) -准素模 Q 的直和。（化归到 M 不含非零可除子模的情形，并用反证法论证。利用 b) 证明：将存在 M 中高度为 1 的元素组成的无穷序列 \((x_k)\) ，使得：1) \(x_k\) 的零化子是 \(\mathrm{A}\pi^{n_k}\) ，其中整数序列 \((n_k)\) 是严格递增的；并且 2）诸 \(\mathrm{A}x_k\) 的和是直和，且这些子模中任意有限个的和是 M 的一个直和因子。然后证明 M 的由元素 \(x_k - \pi^{n_{k+1} - n_k}x_{k+1}\) 生成的子模 N 是纯的，但不是 M 的直因子，论证如习题 7 e)）。
\end{enumerate}

\begin{enumerate}
\def\labelenumi{\arabic{enumi})}
\setcounter{enumi}{8}
\tightlist
\item
  \begin{enumerate}
  \def\labelenumii{\alph{enumii})}
  \tightlist
  \item
    设 M 为一个 \(\pi\) -准素模，并设 \(M_{0}\) 为 M 中由无限高度元素组成的子模。证明 \(\pi\) -准素模 \(M/M_{0}\) 不包含任何具有无限高度的非零元素。
  \end{enumerate}
\end{enumerate}

\begin{enumerate}
\def\labelenumi{\alph{enumi})}
\setcounter{enumi}{1}
\tightlist
\item
  设 \(p\) 是一个素数，且令 \(E\) 是 \(p\) -群 \(E_n = \mathbf{Z} / (p^n)\) ( \(n \geqslant 1\) ) 的直和；令 \(e_n\) 表示 \(1 \mod p^n\) 在 \(E_n\) 中的剩余类。设 \(H\) 是 \(E\) 的由元素 \(e_1 - p^{n-1}e_n\) （对所有整数 \(n \geqslant 2\) ）生成的子模。设 \(M\) 是商模 \(E / H\) 。证明由无限高度元素构成的子模 \(M_0\) （ \(M\) 中的）是子模 \((E_1 + H) / H\) ，它同构于 \(E_1\) ，并证明 \(M_0\) 中不存在在 \(M_0\) 中具有无限高度的非零元素。
\end{enumerate}

¶ 10) a) 设 M 是一个 π-准素模，且没有非零的无限高度元素，并设 N 是 M 的纯子模，由 M 的一个极大拟自由元素族生成（习题 4 和 7，b)）。证明 M/N 是一个可除模（用反证法论证，依次应用习题 8，c)、习题 7，c)、习题 7，b)，最后应用习题 4，c））。

\begin{enumerate}
\def\labelenumi{\alph{enumi})}
\setcounter{enumi}{1}
\item
  考虑 M 上如下定义的可度量化拓扑 \(\mathcal{T}\) ：取子模 \(\pi^n\mathbf{M}\) （对所有整数 \(n\geqslant 0\) ）作为 0 的邻域基（一般拓扑，III，第5页）。那么，M 关于子模 P 的商 M/P 是可除的，当且仅当 P 在拓扑 \(\mathcal{T}\) 下在 M 中处处稠密。
\item
  设 \(P_{1}\) , \(P_{2}\) 是 M 的两个子模，它们是纯的，在拓扑 T 下在 M 中处处稠密，并且是循环子模的直和。证明 \(P_{1}\) 和 \(P_{2}\) 是同构的（注意 \(P_{1}/\pi^{n}P_{1}\) 对所有 n 同构于 \(M/\pi^{n}M\) ）。
\item
  设 \(\hat{\mathbf{M}}\) 是 \(\mathbf{M}\) 关于拓扑 \(\mathcal{T}\) 的完备化（一般拓扑学，III，第24页）。证明 \(x \mapsto \alpha x\) 是从 \(\hat{\mathbf{M}}\) 到 \(\alpha \hat{\mathbf{M}}\) 上的严格态射，对所有 \(\alpha \in \mathbf{A}\) ， \(\pi^n \hat{\mathbf{M}}\) 是开且闭的子模，且是 \(\pi^n \mathbf{M}\) 的闭包（对所有 \(n\) ），并且商模 \(\hat{\mathbf{M}} / \pi^n \hat{\mathbf{M}}\) 与 \(\mathbf{M} / \pi^n \mathbf{M}\) 是同构的。
\item
  设 \(\bar{\mathbf{M}}\) 为 \(\hat{\mathbf{M}}\) 的挠子模；证明 \(\bar{\mathbf{M}}\) 是一个 \(\pi\) -准素模，有 \(\bar{\bar{\mathbf{M}}} = \bar{\mathbf{M}}\) ，并且 \(\mathbf{M}\) 是 \(\bar{\mathbf{M}}\) 的纯子模。
\item
  假设 \(\pi\) -准素模 M 是循环子模的一个族 \((E_{\alpha})\) 的直和。那么 T 是离散的当且仅当 M 的零化子是 0。在这种情况下，令 E 为积模 \(\prod_{\alpha} E_{\alpha}\) ，并令 F 为 E 的由这样的元素组成的子模：这些元素除
\end{enumerate}

可数多个坐标外所有坐标均为零；证明 \(\hat{\mathbf{M}}\) 随后可以（作为非拓扑模）与 \(F\) 的由这样的 \((x_{\alpha})\) 组成的子模等同：对于每个整数 \(n\) ，只有有限多个指标 \(\alpha\) 使得 \(x_{\alpha} \notin \pi^{n} E_{\alpha}\) 。然后证明这四个模 \(\mathbf{M}\) , \(\hat{\mathbf{M}}\) , \(\bar{\mathbf{M}}\) 和 \(F\) 是不同的。由此推出 \(\bar{\mathbf{M}}\) 不是循环模的直和（使用 \(e\) ），并且它不是不可分解模的直和（注意 \(\bar{\mathbf{M}}\) 不含无限高度的元素，并应用习题8，c））。

\begin{enumerate}
\def\labelenumi{\arabic{enumi})}
\setcounter{enumi}{10}
\tightlist
\item
  设 M 为一个 \(\pi\) -准素模，设N是M的一个子模，其非零元素在M中的高均为 \(\leqslant h\) 。
\end{enumerate}

\begin{enumerate}
\def\labelenumi{\alph{enumi})}
\item
  设 P 是 M 的一个纯子模，使得 \(P(\pi) \subset N(\pi)\) 且 \(P(\pi) \neq N(\pi)\) 。设 a 是 \(N(\pi)\) 中不属于 \(P(\pi)\) 的一个元素，并设 \(x_{0}\) 为 \(a + P(\pi)\) 中的一个元素，其在 M 中的高度 k 尽可能大；设 \(x_{0} = \pi^{k} y_{0}\) 。证明和 \(P + Ay_{0}\) 是直和，并且是 M 的一个纯子模。
\item
  设 B 是 M 的一个拟自由子集，使得若 P 是由 B 生成的纯子模，则 \(\mathrm{P}(\pi) \subset \mathrm{N}(\pi)\) 。证明存在 M 的一个拟自由子集 \(C \supset B\) ，使得若 Q 是由 C 生成的 M 的纯子模，则 \(\mathrm{Q}(\pi) = \mathrm{N}(\pi)\) （应用 Zorn 引理，利用 a)）。
\end{enumerate}

\begin{enumerate}
\def\labelenumi{\arabic{enumi})}
\setcounter{enumi}{11}
\tightlist
\item
  \(\pi\) -准素模 M 是循环子模的直和，当且仅当存在 M 中的拟自由子集 H，使得若 N 是由 H 生成的 M 的纯子模，则 \(M(\pi) = N(\pi)\) 。
\end{enumerate}

\begin{enumerate}
\def\labelenumi{\alph{enumi})}
\setcounter{enumi}{1}
\item
  模 M 是循环子模的直和，当且仅当 M 是子模的某个递增序列 \((P_{k})\) 的并，使得每个 \(P_{k}\) 中非零元素在 M 中的高度是有界的。（要看出该条件是充分的，使用 a) 和习题 11 的 b)）。
\item
  由 \(b)\) 推出，若 \(M\) 是循环子模的直和，则 \(M\) 的每个子模都是循环子模的直和。
\item
  设 M 为一个 \(\pi\) -准素模，没有无限高度元素，具有可数生成系（ \(x_{k}\) ）。证明 M 是循环子模的直和（若
\end{enumerate}

\(P_{k}\) 是由 \(x_{i}\) （给定 \(i \leqslant k\) ）生成的子模，则将 b) 应用于递增序列 \((\mathrm{P}_{k})\) ，并利用习题 5 应用于递减序列 \((\mathrm{P}_{k} \cap \pi^{j}\mathrm{M})\) （对固定的 k））。

\begin{enumerate}
\def\labelenumi{\alph{enumi})}
\setcounter{enumi}{4}
\tightlist
\item
  设 M 为一个 \(\pi\) -准素模，具有可数生成系。则 M 是不可分解子模的直和当且仅当 M 中无限高度元素构成的子模 N 是纯子模（利用 d) 和习题 7，c)）。
\end{enumerate}

\begin{enumerate}
\def\labelenumi{\arabic{enumi})}
\setcounter{enumi}{12}
\tightlist
\item
  设 A 为主理想整环，设 \(\pi\) 为 A 的一个不可约元素，并设 M 为 A 上的 \(\pi\) -准素模。设 I 为满足如下条件的序数 \(\alpha\) （集合论，III，第 225 页，习题14）组成的集：序数之集 \(< \alpha\) 的基数至多等于 M 的基数。由超限归纳法，定义 M 的子模 \(\mathbf{M}^{(\alpha)}\) ——对所有 \(\alpha \in I\) ——如下：取 \(\mathbf{M}^{(0)} = \mathbf{M}\) ；取 \(\mathbf{M}^{(\alpha + 1)} = \pi \mathbf{M}^{(\alpha)}\) ；并且如果 \(\alpha\) 是极限序数，取 \(\mathbf{M}^{(\alpha)}\) 为所有 \(\mathbf{M}^{(\beta)}\) （给定 \(\beta < \alpha\) ）的交集。
\end{enumerate}

\begin{enumerate}
\def\labelenumi{\alph{enumi})}
\item
  证明存在一个最小序数 \(\tau \in I\) 使得 \(M^{(\tau + 1)} = M^{(\tau)}\) （M 的末端序数）。证明 \(M^{(\tau)}\) 在 M 中有补，并且是若干同构于 \(U_{\pi}\) 的子模的直和（习题 3）。称 M 是约化的，如果 \(M^{(\tau)} = 0\) 。
\item
  设 M 为约化 \(\pi\) -模。对所有 \(x \in \mathbf{M}\) ，证明存在一个最大序数 \(\gamma(x) \leqslant \tau\) 使得 \(x \in \mathbf{M}^{(\gamma(x))}\) （特别地， \(\gamma(0) = \tau\) ）。推广习题 8 中的定义，我们称 \(\gamma(x)\) 为 \(x\) 在 M 中的高度。证明若 \(\gamma(x) < \gamma(y)\) ，则 \(\gamma(x + y) = \gamma(x)\) ，且若 \(\gamma(x) = \gamma(y)\) ，则 \(\gamma(x + y) \geqslant \gamma(x)\) 。我们称 \(x\) 关于 M 的子模 N 是适当的，如果对所有 \(y \in \mathbf{N}\) ，高度 \(\gamma(x + y)\) 至多等于 \(\gamma(x)\) ；证明此时 \(\gamma(x + y) = \inf(\gamma(x), \gamma(y))\) 对所有 \(y \in \mathbf{N}\) 成立。
\item
  设 \(\alpha < \tau\) 为一个序数，并设 \(\tilde{\mathbf{M}}^{(\alpha)}\) 为 \(\mathbf{M}^{(\alpha)}\) 的子模，由那些 \(x\) （使得 \(\pi x = \pi y\) 对某个 \(y \in \mathbf{M}^{(\alpha + 1)}\) ）组成。对于那些具有此性质的 \(y \in \mathbf{M}^{(\alpha + 1)}\) ，形如 \(u = x - y\) 的元素构成 模 \(\mathbf{M}(\pi) \cap \mathbf{M}^{(\alpha + 1)}\) 在 \(\mathbf{M}(\pi) \cap \mathbf{M}^{(\alpha)}\) 中的一个类；如果 \(\varphi\) 是从 \(\mathbf{M}(\pi) \cap \mathbf{M}^{(\alpha)}\) 到该模关于 \(\mathbf{M}(\pi) \cap \mathbf{M}^{(\alpha + 1)}\) 的商上的典范同态，则证明映射 \(\psi\) ——它将 \(x\) 映为 \(\varphi(u)\) ——是一个核为 \(\mathbf{M}^{(\alpha + 1)}\) 的同态，并且因此 \(\tilde{\mathbf{M}}^{(\alpha)} / \mathbf{M}^{(\alpha + 1)}\) 同构于 \((\mathbf{M}(\pi) \cap \mathbf{M}^{(\alpha)}) / (\mathbf{M}(\pi) \cap \mathbf{M}^{(\alpha + 1)})\) 。
\item
  设 N 是 M 的一个子模，并设 \(x \in \tilde{\mathbf{M}}^{(\alpha)}\) 是高度为 \(\alpha\) 的元素；证明 x 关于 N 是适当的，当且仅当 \(\mathbf{M}(\pi) \cap \mathbf{M}^{(\alpha)}\) 中对应的元素 u = x - y （参见 c)）关于 N 是适当的。存在这样的元素 x，当且仅当 \(\mathbf{N} \cap \tilde{\mathbf{M}}^{(\alpha)}\) 在同态 \(\psi\) 下的像不同于 \((\mathbf{M}(\pi) \cap \mathbf{M}^{(\alpha)}) / (\mathbf{M}(\pi) \cap \mathbf{M}^{(\alpha + 1)})\) 。
\end{enumerate}

\begin{enumerate}
\def\labelenumi{\arabic{enumi})}
\setcounter{enumi}{13}
\tightlist
\item
  设 M 为约化 \(\pi\) -准素模。在习题 13 的记号中，对于每个序数 \(\alpha < \tau\) ，我们可以把 \((\mathbf{M}(\pi) \cap \mathbf{M}^{(\alpha)}) / (\mathbf{M}(\pi) \cap \mathbf{M}^{(\alpha + 1)})\) 看作域 \(F_{\pi} = A / A\pi\) 上的向量空间；设 \(c(\alpha)\) 表示这个向量空间的维数，并称 \(c(\alpha)\) 为序数 \(\alpha\) 在模 M 中的重数（参见 VII，第 24 页）。
\end{enumerate}

\begin{enumerate}
\def\labelenumi{\alph{enumi})}
\item
  证明：若 M 和 N 是两个 \(\pi\) -准素模，且每个都是循环子模的直和（这意味着 \(\tau \leqslant \omega\) ，即第一个无限序数），则 M 与 N 同构当且仅当它们具有相同的末端序数 \(\tau\) ，并且每个整数 \(n < \tau\) 的重数在 M 中与在 N 中相同。
\item
  给出两个没有无限高度元素的 \(\pi\) -准素模的例子，它们不同构，但其中每个整数 n 具有相同的重数（参见习题 10，f）。c) 设 M 和 N 是两个约化 \(\pi\) -准素模，每个都具有可数无限生成系。证明如果 M 和 N 具有相同的末端序数 \(\tau\) ，并且每个序数 \(\alpha < \tau\) 在 M 和 N 中具有相同的重数，则 M 与 N 同构（« Ulm–Zippin 定理 »）。（给定两个子模 \(S \subset M\) 和 \(T \subset N\) ，一个同构 \(h\) （从 \(S\) 到 \(T\) 上）称为指标同构，如果对所有 \(x \in S\) ， \(x\) 在 \(M\) 中的高度等于 \(h(x)\) 在 \(N\) 中的高度。通过归纳法，并交替交换 \(M\) 和 \(N\) 所扮演的角色，定义一个从 \(M\) 到 \(N\) 上的指标同构，使得问题归结为如下问题：给定两个有限生成的子模 \(S \subset M\) 和 \(T \subset N\) ，一个指标同构 \(h\) （从 \(S\) 到 \(T\) 上），以及一个元素 \(x \in M \cap \complement S\) ，将 \(h\) 延拓为一个从 \(S + Ax\) 到 \(N\) 的一个包含 \(T\) 的子模上的指标同构。证明这可以归结为 \(\pi x \notin S\) 的情形；利用习题 5，选取一个元素 \(w_1\) ，它属于 \((S + Ax) \cap \complement S\) ，具有最大可能的高度 \(\alpha\) （在 \(M\) 中）；然后选取一个元素 \(w_2\) ，它属于 \(S \cap M^{(\alpha)}\) ，使得 \(\pi(w_1 + w_2)\) 在 \(M\) 中的高度是最大可能的；证明 \(w = w_1 + w_2\) 关于 \(S\) 是适当的（习题13，b)）。令 \(z = h(\pi w) \in T\) ；证明，如果存在一个元素 \(t\) （属于 \(N\) ，高度为 \(\alpha\) （在 \(N\) 中），关于 \(T\) 是适当的，且满足 \(\pi t = z\) ），那么 \(\bar{h}\) （ \(h\) 的一个所需形式的扩张）可以通过令 \(\bar{h}(w) = t\) 来得到。剩下的问题是证明元素 \(t\) 的存在性。若
\end{enumerate}

\[
\gamma (z) = \gamma (\pi w) = \alpha + 1 <   \tau ,
\]

即证明存在这样的元素 \(u\) ：它属于 \(\mathbf{N}^{(\alpha)} \cap \complement \mathbf{T}\) ，且 \(z = \pi u\) ，并且 \(t\) 可以取为任意这样的元素。另一方面，如果 \(\gamma(z) = \gamma(\pi w) > \alpha + 1\) ，或者如果 \(\gamma(z) = \gamma(\pi w) = \alpha + 1 = \tau\) ，则利用习题 13， \(d\) ) 来证明元素 \(t\) 的存在性，注意到（在习题 13 的记号下）在 \(\psi\) 之下， \(S \cap \tilde{\mathbf{M}}^{(\alpha)}\) 在 \((\mathbf{M}(\pi) \cap \mathbf{M}^{(\alpha)}) / (\mathbf{M}(\pi) \cap \mathbf{M}^{(\alpha + 1)})\) 中的像是域 \(F_{\pi}\) 上的有限维空间）

* d) 由 c) 得出习题 12 d) 的一个新证明（同时使用第七章第14页的定理1）。 *

\begin{enumerate}
\def\labelenumi{\arabic{enumi})}
\setcounter{enumi}{14}
\tightlist
\item
  \begin{enumerate}
  \def\labelenumii{\alph{enumii})}
  \tightlist
  \item
    设 A 为主理想整环，设 M 为可除 A-模（习题 3），并设 N 为挠 A-模。证明 \(\mathbf{M} \otimes_{\mathbf{A}} \mathbf{N} = 0\) 。
  \end{enumerate}
\end{enumerate}

\begin{enumerate}
\def\labelenumi{\alph{enumi})}
\setcounter{enumi}{1}
\tightlist
\item
  设 M 和 N 为两个没有无限高度元素的 \(\pi\) -准素模；证明 \(M \otimes_{A} N\) 是循环子模的直和（利用 a) 以及习题 10，a)）。
\end{enumerate}

\BourbakiExerciseGroup

\begin{enumerate}
\def\labelenumi{\arabic{enumi})}
\item
  设 A 是一个交换环。证明：如果每个自由 A-模的每个子模都是自由的，那么 A 的每个理想都是主理想，并且 A 是整环（从而是一个主理想整环）。
\item
  设 M 是主理想整环 A 上的一个模；假设 M 是循环子模的直和。证明 M 的每个子模 N 都是循环子模的直和（若 T 是 N 的挠子模，首先利用定理 1 证明 N/T 是自由模，并由此推出 T 在 N 中有补；然后利用第七章第 57 页习题 12 c) 证明 T 是循环子模的直和）。
\item
  设 \((p_n)\) 是一个严格递增的素数序列，使得 \(p_n\) 不整除下列数中的任何一个
\end{enumerate}

\[
a + b \left(1 + p _ {1} + p _ {1} p _ {2} + \dots + p _ {1} p _ {2} \dots p _ {n - 1}\right) \quad \text { for } \quad | a | \leqslant n \quad \text { and } \quad | b | \leqslant n.
\]

在有理平面 \(\mathbf{Q}^2\) 中，考虑 \(\mathbf{Z}\) -模 \(\mathbf{M}\) ，它由如下定义的序列 \((x_n)\) 生成：

\[
x _ {0} = (1, 0), \quad x _ {1} = (0, 1), \quad x _ {n + 1} = p _ {n} ^ {- 1} \left(x _ {0} + x _ {n}\right) \quad \text { for } \quad n \geqslant 1.
\]

证明 M 的秩为 2，并且 M 的每个秩 1 子模都是循环的（注意 \(x_{0}\) 和 \(x_{n}\) 构成 M 的由诸 \(x_{i}\) （满足 \(i \leqslant n\) ）生成的子模的一组基）。

\begin{enumerate}
\def\labelenumi{\arabic{enumi})}
\setcounter{enumi}{3}
\tightlist
\item
  设 \(E\) 是主理想整环 \(A\) 上的无挠模。
\end{enumerate}

\begin{enumerate}
\def\labelenumi{\alph{enumi})}
\item
  证明：若F是E的纯子模（VII，第55页，习题7），则E/F是无挠的；反之亦然。
\item
  设 F 为 E 的由一个极大拟自由族（VII，第 55 页，习题 4）生成的子模；则 F 是 E 的纯的自由子模。证明：对所有 \(\dot{x} \in \mathrm{E} / \mathrm{F}\) ，存在一个属于 A 的不可逆元 \(\alpha\) 以及一个元素 \(\dot{y} \in \mathrm{E} / \mathrm{F}\) ，使得 \(\dot{x} = \alpha \dot{y}\) ，但 E/F 不一定是可除模（参见习题 3）。
\end{enumerate}

\begin{enumerate}
\def\labelenumi{\arabic{enumi})}
\setcounter{enumi}{4}
\tightlist
\item
  设 A 是主理想整环，且具有唯一的极大理想 \(\mathbf{A}\pi\) （参见 VII，第 48 页，习题 4）。设 E 是无挠 A-模，M 是 E 的自由子模，使得 E/M 是无挠、可除且秩为 1 的。设 \((e_{\alpha})\) 是 M 的一组基，设 \(\dot{u}\) 是 E/M 的一个非零元素，且 \(u\) 是陪集 \(\dot{u}\) 的一个元素；对每个整数 \(n > 0\) ，设 \(u_{n}\) 是 E 的一个元素，使得 \(u \equiv \pi^{n} u_{n} (\bmod M)\) 。令 \(u = \pi^{n} u_{n} + \sum_{\alpha} \lambda_{n\alpha} e_{\alpha}\) 。
\end{enumerate}

\begin{enumerate}
\def\labelenumi{\alph{enumi})}
\item
  证明 M 由 M 和诸 \(u_{n}\) 生成，并且 \(\lambda_{n\alpha} - \lambda_{n + 1,\alpha} \in \mathbf{A}\pi^n\) 对所有 \(\alpha\) 并且对每个整数 \(n > 0\) 成立（参见习题 4，a)）。
\item
  进一步假设 A 关于拓扑 T 是完备的，T 由取理想 \(A\pi^{n}\) 作为 A 中 0 的开邻域基而定义（一般拓扑，III，第 5 页）；令 \(\lambda_{\alpha} = \lim_{n \to \infty} \lambda_{n\alpha}\) 。证明：如果 \(\mu\) 和 \(\mu_{\alpha}\) 使得 \(\mu u - \sum_{\alpha} \mu_{\alpha} e_{\alpha}\) 属于所有
\end{enumerate}

子模 \(\pi^n\mathrm{E}\) ，则必然有 \(\mu_{\alpha} = \lambda_{\alpha}\mu\) 对所有 \(\alpha\) 。

\begin{enumerate}
\def\labelenumi{\arabic{enumi})}
\setcounter{enumi}{5}
\tightlist
\item
  设 \(\mathbf{A}\) 是一个主理想整环，具有唯一的极大理想 \(\mathbf{A}\pi\) ，并且关于习题 5 中定义的拓扑 \(\mathcal{T}\) 是完备的。
\end{enumerate}

\begin{enumerate}
\def\labelenumi{\alph{enumi})}
\item
  证明：任意秩有限的、无挠的 A-模 E，若不包含非零可除子模，则 E 是自由的（考虑由一个极大拟自由族生成的 E 的子模 M；若 \(M \neq E\) 则借助习题 4，b) 和 5，b) 证明 E 包含非零可除子模）。将此结果推广到 E 具有可数秩的情形（用归纳法进行）。
\item
  在自由模 \(E = A^{(N)}\) （A 的可数无穷多个副本的直和）中，设 \((a_{n})_{n \geq 0}\) 为典范基，并令 \(e_{n} = a_{n-1} - \pi a_{n}\) 对每个整数 \(n \geq 1\) 。证明由诸 \(e_{n}\) 生成的 E 的子模 M 是纯的，并且诸 \(e_{n}\) 构成一个极大的拟自由族，但 E/M 是可除模，因此 M 在 E 中没有补。
\end{enumerate}

\begin{enumerate}
\def\labelenumi{\arabic{enumi})}
\setcounter{enumi}{6}
\tightlist
\item
  设 \(p\) 是一个素数，并设 \(\mathbf{A}\) 表示形如 \(r / s\) 的有理数组成的环，其中 \(s\) 与 \(p\) 互素（VII，第 48 页，习题 4）。
\end{enumerate}

设 \(u = (1,0)\) 和 \(v = (0,1)\) 是向量空间 \(E = Q^2\) （ \(Q\) 上的）的典范基的元素。递归地定义一个元素序列 \(u_n\) （ \(E\) 中的）： \(u_0 = u\) ，且 \(u_{n}=pu_{n+1}+v\) 对 \(n\geqslant0\) 。设 F 为由 v 和诸 \(u_{n}\) 生成的 A-子模；证明 F 不是自由模，但不含非零可除子模（参见习题 6，a)）。

\begin{enumerate}
\def\labelenumi{\arabic{enumi})}
\setcounter{enumi}{7}
\tightlist
\item
  设 A 是一个非域的主理想整环。
\end{enumerate}

\begin{enumerate}
\def\labelenumi{\alph{enumi})}
\item
  证明乘积 A-模 \(A^{N}\) 没有可数生成系。（若 A 可数，考虑 \(\operatorname{Card}(A^{N})\) ；否则注意到，若对所有 \(x \in A\) 我们令 \(v_{x}\) 表示元素 \((x^{n})_{n \geq 0}\) （属于 \(A^{N}\) ），则诸 \(v_{x}\) 构成自由系）。
\item
  设 \(\pi\) 为 A 的一个不可约元素。令 S 表示 \(\mathbf{A}^{\mathbf{N}}\) 的子模，它由这样的序列 \(z = (z_{n})_{n\in \mathbb{N}}\) 组成：存在一个整数序列 \((k(n))_{n\in \mathbb{N}}\) ，它趋于 \(+\infty\) （依赖于 \(z\) ），使得 \(z_{n}\in \mathbf{A}\pi^{k(n)}\) 对所有 \(n\) 成立。证明 \((\mathbf{A} / \pi)\) -向量空间 \(\mathrm{S} / \pi \mathrm{S}\) 有一个可数基。
\item
  由 a) 和 b) 推出， \(A^{N}\) 不是自由 A-模。（否则，S 将是一个具有可数基的自由 A-模；但 \((x_{n})\mapsto(x_{n}\pi^{n})\) 是一个从 \(A^{N}\) 到 S 中的单射 A-线性映射。）
\end{enumerate}

\begin{enumerate}
\def\labelenumi{\arabic{enumi})}
\setcounter{enumi}{8}
\tightlist
\item
  设 \(\mathbf{A}\) 为主理想整环，且设 \(\mathbf{M}\) 为 \(\mathbf{A}\) -模 \(\mathbf{A}^{\mathbf{N}}\) 的对偶模。设 \(f: \mathbf{M} \to \mathbf{A}^{\mathbf{N}}\) 为典范单射 \(\mathbf{A}^{(\mathbf{N})} \to \mathbf{A}^{\mathbf{N}}\) 的转置映射（其中 \(\mathbf{A}^{\mathbf{N}}\) 与 \(\mathbf{A}^{(\mathbf{N})}\) 的对偶等同）。
\end{enumerate}

\begin{enumerate}
\def\labelenumi{\alph{enumi})}
\item
  设 \(\pi\) 为 A 的一个不可约元素，且令 \(\varphi: \mathbf{A}^{\mathbf{N}} / \mathbf{A}^{(\mathbf{N})} \to \mathbf{A}\) 是一个线性形式；证明如果 \(u \in \mathbf{A}^{\mathbf{N}} / \mathbf{A}^{(\mathbf{N})}\) 是某个元素 \((u_n) \in \mathbf{A}^{\mathbf{N}}\) （满足 \(u_n \in (\pi^n)\) 对所有 \(n\) ）的类，则 \(\varphi(u) = 0\) 。如果 A 包含两个非相伴的不可约元素，则推出 \(f\) 是单射的。（利用上述结论和贝祖恒等式，证明 \(\mathbf{A}^{\mathbf{N}} / \mathbf{A}^{(\mathbf{N})}\) 的对偶为零。）
\item
  给 A 赋予 VII 第 60 页习题 5 中所定义的拓扑。证明：若 \(f(\mathbf{M})\) 不包含于 \(\mathbf{A}^{(\mathbf{N})}\) ，则 A 是完备的。（将问题归结为证明：若 \(\varphi : \mathbf{A}^{\mathbf{N}} \to \mathbf{A}\) 是一个线性形式，使得 \(\varphi(e_n) \neq 0\) 对每个元素 \(e_n\) （属于 \(\mathbf{A}^{(\mathbf{N})}\) 的典范基）成立，则 \(\operatorname{Im}(\varphi)\) 包含 A 中任意一个趋于 0 足够快的元素级数的和。）
\item
  如果 A 包含两个非相伴的不可约元素，则从 \(\mathbf{A}^{(N)}\) （相应地 \(A^{N}\) ）到其双重对偶的典范映射是双射（利用 a) 和 b)）。
\end{enumerate}

\begin{enumerate}
\def\labelenumi{\arabic{enumi})}
\setcounter{enumi}{9}
\item
  设 M 为 p 进整数环 \(Z_{p}\) 的一个非零 Z-子模，且是 Z-模 \(Z_{p}\) （VII，第 55 页，习题 7）的纯 Z-子模。证明 \(M + pZ_{p} = Z_{p}\) ；由此推出 M 是不可分解的（注意 \(Z_{p}/pZ_{p}\) 同构于 Z/pZ）。给出 \(Z_{p}\) 的一个可分解 Z-子模的例子（观察到 \(Z_{p}\) 具有连续统的基数）。
\item
  设 A 和 B 是两个无限的有理素数集合，它们没有公共元素，且都不包含数 5。设 \((e_i)_{1 \leqslant i \leqslant 4}\) 是 Q-向量空间 \(\mathbf{Q}^4\) 的典范基。
\end{enumerate}

\begin{enumerate}
\def\labelenumi{\alph{enumi})}
\item
  设 M 为 \(Q^{4}\) 的由元素 \(e_{1}/m\) （给定 \(m \in A\) ）生成的 Z-子模；设 N 为 \(Q^{4}\) 的由元素 \(e_{2}/m\) （其中 \(m \in A\) ）、 \(e_{3}/n\) （其中 \(n \in B\) ）以及 \((e_{2} + e_{3})/5\) 生成的 Z-子模，则 \(M \cap N = 0\) 。令 \(e_{1}' = 8e_{1} + 3e_{2}\) 和 \(e_{2}' = 5e_{1} + 2e_{2}\) 。设 \(M'\) 为 \(Q^{4}\) 的由元素 \(e_{1}'/m\) （其中 \(m \in A\) ）生成的 Z-子模；设 \(N'\) 为 \(Q^{4}\) 的由元素 \(e_{2}'/m\) （其中 \(m \in A\) ）、 \(e_{4}/n\) （其中 \(n \in B\) ）以及 \((3e_{2}' + e_{3})/5\) 生成的 Z-子模。证明 \(M' \cap N' = 0\) ， \(M \oplus N = M' \oplus N'\) ；M 与 \(M'\) 同构且不可分解；并且 N 与 N' 不可分解但不同构（比较 VII，第 69 页，习题 23）。
\item
  设 M 为 \(Q^{4}\) 的由元素 \(e_{1}/m\) （其中 \(m \in A\) ）、 \(e_{2}/n\) （其中 \(n \in B\) ）以及 \((e_{1} + e_{2})/5\) 生成的 Z-子模；设 N 为 \(Q^{4}\) 的由元素 \(e_{3}/m\) （其中 \(m \in A\) ）、 \(e_{4}/n\) （其中 \(n \in B\) ）以及 \((e_{3} + e_{4})/5\) 生成的 Z-子模。令 \(e_{1}' = 2e_{1} + e_{3}\) , \(e_{2}' = e_{2} + 3e_{4}\) , \(e_{3}' = 17e_{1} + 9e_{3}\) ，以及 \(e_{4}' = e_{2} + 2e_{4}\) 。设 \(M'\) 为 \(Q^{4}\) 的由元素 \(e_{1}'/m\) （给定 \(m \in A\) ）、 \(e_{2}'/n\) （给定 \(n \in B\) ）以及 \((e_{1}' + 2e_{2}')/5\) 生成的 Z-子模；设 \(N'\) 为 \(Q^{4}\) 的由元素 \(e_{3}'/m\) （给定 \(m \in A\) ）、 \(e_{4}'/n\) （给定 \(n \in B\) ）以及 \((e_{3}' + 2e_{4}')/5\) 生成的 Z-子模。证明 \(M \oplus N = M' \oplus N'\) ，M 与 N 同构，且 \(M'\) 同构于 \(N'\) ，但 \(M'\) 与 M 不同构。
\item
  设 M 为 \(Q^{4}\) 的由 \(e_{1}/5^{n}\) ( \(n \geqslant 0\) ) 生成的 Z-子模；设 N 为 \(Q^{4}\) 的由 \(e_{2}/5^{n}\) , \(e_{3}/2^{n}\) , \(e_{4}/3^{n}\) ( \(n \geqslant 0\) ), \((e_{2} + e_{3})/3\) 和 \((e_{2} + e_{4})/2\) 生成的子模。令 \(e_{1}' = 3e_{1} - e_{2}\) 和 \(e_{2}' = 2e_{1} - e_{2}\) 。设 \(M'\) 为 \(Q^{4}\) 的由元素 \(e_{1}'/5^{n}\) , \(e_{3}/2^{n}\) ( \(n \geqslant 0\) ）和 \((e_{1}' - e_{3})/8\) 生成的 Z-子模；设 \(N'\) 为 \(Q^{4}\) 的由元素 \(e_{2}'/5^{n}\) , \(e_{4}/3^{n}\) ( \(n \geqslant 0\) ）和 \((e_{2}' - e_{4})/2\) 生成的 Z-子模。证明 \(M \oplus N = M' \oplus N'\) ， \(M'\) 和 \(N'\) 是秩为 2 的不可分解模，且 M 的秩为 1。
\end{enumerate}
